% =============================================================================
% Assignments/CS301/09-efficient-sorting-coding-assignment.tex
% Coding Assignment 9: Merge, Quick, and Shell Sort
% (Week 9, syllabi/CS301.md Week 9 row: P9 quick sort, P10 shell sort,
% P11 merge sort)
%
% Student-facing problem statement. Instructor-only files live in
% Assignments/CS301/09-efficient-sorting-coding/:
%   reference_solution.c  (NEVER published, gitignored)
%   test_hidden.c         (NEVER published, gitignored)
%   autograder.sh         (NEVER published, gitignored)
%   sorting.h, starter_sorting.c, test_visible.c  (provided to students)
%
% Fresh instances: the deck's own examples (merge [3,8,15]+[2,9,11];
% partition [7,2,9,4,3,8,5]; quick on [1..6]; radix [170,45,75,90,802,24];
% shell [7,3,9,1,8,2,6,4]; heap [50,30,40,10,20,35,25]) and the written
% Assignment 9 instances (merge [2,6,11]+[1,4,9,10]; partition
% [8,3,7,4,2,9,5]; quick [2,4,6,8,10]) are NOT reused. Visible tests use
% [1,4,7 | 2,5,8], [5,1,4,2,8,0,3], [3,1,2,1,3,2], and inversion arrays.
%
% Compile with (macOS/Linux): cd Assignments/CS301 &&
%   TEXINPUTS=../../Preambles:$TEXINPUTS xelatex -interaction=nonstopmode 09-efficient-sorting-coding-assignment.tex (3 passes)
% =============================================================================

\documentclass[11pt]{article}
\usepackage[margin=1in]{geometry}
% =============================================================================
% Preambles/header.tex — shared LaTeX/Beamer preamble for this project
%
% Use from a Beamer lecture by prepending:
%     \input{header}
% and compiling with TEXINPUTS=../Preambles:$TEXINPUTS (the /compile-latex
% skill does this automatically).
%
% PALETTE CONTRACT. Color names defined here MUST match the names in
% Quarto/theme-template.scss so Beamer and Quarto renderings use the same
% palette. The scripts/check-palette-sync.sh script enforces this.
%
% To customize: edit the HEX values below AND the matching $variable in
% Quarto/theme-template.scss, then run scripts/check-palette-sync.sh.
% =============================================================================

% -----------------------------------------------------------------------------
% Engine detection + encoding
%
% Our /compile-latex skill uses XeLaTeX, where inputenc is unnecessary and
% can produce encoding warnings. Only load inputenc under pdfTeX.
% -----------------------------------------------------------------------------
\usepackage{iftex}
\ifPDFTeX
  \usepackage[utf8]{inputenc}
\fi

\usepackage{amsmath, amssymb, amsthm}
\usepackage{booktabs}
\usepackage{graphicx}

% -----------------------------------------------------------------------------
% PALETTE — mirrors Quarto/theme-template.scss variable names.
% Load xcolor and define colors BEFORE anything that references them
% (notably hyperref, hypersetup, and the Beamer color assignments).
% -----------------------------------------------------------------------------
\usepackage{xcolor}

\definecolor{primary-blue}{HTML}{012169}      % headings, accents
\definecolor{primary-gold}{HTML}{B9975B}      % emphasis, borders
\definecolor{highlight-yellow}{HTML}{F2A900}  % markers, alerts
\definecolor{light-bg}{HTML}{E8EDF5}          % title / section backgrounds
\definecolor{jet}{HTML}{1A1A1A}               % body text

% Semantic colors (used in diagrams and annotations)
\definecolor{positive}{HTML}{15803D}          % good / observed / identified
\definecolor{negative}{HTML}{B91C1C}          % bad / problematic / bias
\definecolor{neutral}{HTML}{525252}           % reference / context

% Highlight accents (match .hi-* classes in the SCSS)
\definecolor{hi-slate}{HTML}{314F4F}
\definecolor{hi-green}{HTML}{2E8B57}
\definecolor{hi-red}{HTML}{C41E3A}

% Semantic aliases so skill templates and snippets can write
% \textcolor{accent}{...} without hard-coding "primary-blue".
\colorlet{accent}{primary-blue}
\colorlet{accent2}{primary-gold}

% -----------------------------------------------------------------------------
% Hyperref — loaded AFTER xcolor + palette so link colors resolve.
% -----------------------------------------------------------------------------
\usepackage{hyperref}
\hypersetup{colorlinks=true, linkcolor=primary-blue, urlcolor=primary-blue,
            citecolor=primary-blue, pdfborder={0 0 0}}

% -----------------------------------------------------------------------------
% Course code — set once per deck (after \input{header}, before \begin{document})
% via \coursecode{CS401}. Shown in the footer on every slide so a deck's course
% is identifiable without opening the title page. Empty by default.
% -----------------------------------------------------------------------------
\makeatletter
\newcommand{\coursecode}[1]{\gdef\@coursecodetext{#1}}
\gdef\@coursecodetext{}
\makeatother

% -----------------------------------------------------------------------------
% Beamer theme assignments (only apply under Beamer).
%
% We detect Beamer via \@ifclassloaded{beamer}, which is the documented,
% non-internal way. This must be wrapped in \makeatletter/\makeatother
% because \@ifclassloaded contains an @-catcode character.
% -----------------------------------------------------------------------------
\makeatletter
\@ifclassloaded{beamer}{%
  \usefonttheme{professionalfonts}
  \setbeamercolor{structure}{fg=primary-blue}
  \setbeamercolor{frametitle}{fg=primary-blue}
  \setbeamercolor{title}{fg=primary-blue}
  \setbeamercolor{subtitle}{fg=primary-gold}
  \setbeamercolor{author}{fg=primary-blue}
  \setbeamercolor{institute}{fg=neutral}
  \setbeamercolor{date}{fg=primary-gold}
  \setbeamercolor{itemize item}{fg=primary-blue}
  \setbeamercolor{itemize subitem}{fg=primary-gold}
  \setbeamercolor{alerted text}{fg=primary-gold}
  \setbeamercolor{block title}{fg=white, bg=primary-blue}
  \setbeamercolor{block body}{bg=light-bg}
  % Semantic block variants: block=definition/concept (blue), exampleblock=
  % worked example (green/positive), alertblock=key takeaway or caution (gold).
  % Keep to <=2 colored boxes per slide across ALL three variants (INV-7).
  \setbeamercolor{block title example}{fg=white, bg=positive}
  \setbeamercolor{block body example}{bg=light-bg}
  \setbeamercolor{block title alerted}{fg=white, bg=primary-gold}
  \setbeamercolor{block body alerted}{bg=light-bg}

  % Minimal footer: course code (left) + page X / N (right), no navigation clutter
  \setbeamertemplate{navigation symbols}{}
  \setbeamertemplate{footline}{%
    \color{neutral}\footnotesize\hspace{0.7em}\@coursecodetext\hfill
    \insertframenumber/\inserttotalframenumber\hspace{0.7em}\vspace{0.3em}}
}{}
\makeatother

% -----------------------------------------------------------------------------
% TikZ setup — libraries the snippet gallery uses
% -----------------------------------------------------------------------------
\usepackage{tikz}
\usetikzlibrary{arrows.meta, positioning, calc, decorations.pathreplacing,
                fit, shapes.geometric, backgrounds}

% Shared TikZ styles — snippets and hand-written diagrams can pick these up
% via `\begin{tikzpicture}[every node/.style={...}]` or by naming specific
% styles (e.g., `\node[dag-node] (X) at (0,0) {X};`).
\tikzset{
  % Node styles for causal DAGs and similar
  dag-node/.style = {
    draw=primary-blue, fill=light-bg, rounded corners=2pt,
    minimum width=1.2cm, minimum height=0.9cm, align=center, font=\small
  },
  decision-node/.style = {
    draw=primary-gold, fill=white, diamond, aspect=1.8,
    minimum width=1.8cm, minimum height=1.2cm, align=center, font=\small,
    inner sep=1pt
  },
  % Edge styles — observed vs counterfactual
  observed-edge/.style = {-{Latex[length=2.5mm]}, thick, draw=jet},
  counterfactual-edge/.style = {-{Latex[length=2.5mm]}, thick, draw=neutral, dashed},
  confound-edge/.style = {-{Latex[length=2.5mm]}, thick, draw=negative, dashed},
  % Data-point styles for plots
  observed-dot/.style = {fill=primary-blue, draw=primary-blue, circle, inner sep=1.2pt},
  counterfactual-dot/.style = {fill=white, draw=neutral, circle, inner sep=1.2pt}
}

% -----------------------------------------------------------------------------
% Convenience macros
% -----------------------------------------------------------------------------
\newcommand{\muted}[1]{\textcolor{neutral}{#1}}
\newcommand{\key}[1]{\textcolor{primary-gold}{\textbf{#1}}}
\newcommand{\good}[1]{\textcolor{positive}{#1}}
\newcommand{\bad}[1]{\textcolor{negative}{#1}}

% Standout frame macro: full-bleed dark background for section transitions.
% Beamer-only (uses \begin{frame}/\setbeamercolor/\usebeamerfont) -- do not
% call from an article-class file (e.g. Notes/<CODE>/*-notes.tex). \muted,
% \key, \good, \bad, and \coursecode above are plain xcolor/gdef and are
% safe to use from either class; verified when Notes/article-class support
% was added.
% Expand: \transitionslide{Your transition title}
\newcommand{\transitionslide}[1]{%
  \begingroup
  \setbeamercolor{background canvas}{bg=primary-blue}
  \begin{frame}[plain]
    \centering
    \vfill
    {\usebeamerfont{title}\color{white}\Large #1\par}
    \vfill
  \end{frame}
  \endgroup
}

% Section divider frame (Beamer-only), an alternative to \transitionslide with a
% darker two-line style: near-black (jet) background, a small white label line
% above a larger gold title, and a thin gold rule along the bottom edge.
% Expand: \sectiondivider{Section 2}{Pointers}
\newcommand{\sectiondivider}[2]{%
  \begingroup
  \setbeamercolor{background canvas}{bg=jet}
  \begin{frame}[plain]
    \vspace*{3.0cm}
    {\color{white}\ttfamily\large #1\par}
    \vspace{0.5cm}
    {\color{primary-gold}\LARGE #2\par}
    \begin{tikzpicture}[remember picture, overlay]
      \draw[primary-gold, line width=2.5pt]
        ($(current page.south west)+(0,0.1)$) -- ($(current page.south east)+(0,0.1)$);
    \end{tikzpicture}
  \end{frame}
  \endgroup
}

% -----------------------------------------------------------------------------
% Channel watermark — "Concepts That Click" (YouTube).
%
% Article-class files (CompetitiveExam Q/A sets, Notes, Assignments) get a
% light diagonal draft watermark on every page plus a centered footer naming
% the channel. Beamer decks get a subtle TikZ background watermark instead
% (the footline already carries the course code). Centralized here so every
% MCQ PDF is branded without per-file edits.
% -----------------------------------------------------------------------------
\makeatletter
\@ifclassloaded{article}{%
  \usepackage{draftwatermark}
  \SetWatermarkText{Concepts That Click}
  \SetWatermarkScale{1}
  \SetWatermarkFontSize{2cm}
  \SetWatermarkLightness{0.86}
  \SetWatermarkAngle{45}
  \usepackage{fancyhdr}
  \pagestyle{fancy}
  \fancyhf{}
  \fancyfoot[C]{\footnotesize\textcolor{neutral}{Concepts That Click~|~YouTube: \href{https://www.youtube.com/@concepts.thatclick}{@concepts.thatclick}}}
  \renewcommand{\headrulewidth}{0pt}
  \renewcommand{\footrulewidth}{0pt}
  \fancypagestyle{plain}{%
    \fancyhf{}%
    \fancyfoot[C]{\footnotesize\textcolor{neutral}{Concepts That Click~|~YouTube: \href{https://www.youtube.com/@concepts.thatclick}{@concepts.thatclick}}}%
    \renewcommand{\headrulewidth}{0pt}%
    \renewcommand{\footrulewidth}{0pt}%
  }
}{}
\@ifclassloaded{beamer}{%
  \setbeamertemplate{background}{%
    \begin{tikzpicture}[remember picture, overlay]
      \node[rotate=45, opacity=0.055, align=center]
        at (current page.center)
        {\Huge\textcolor{neutral}{Concepts That Click}};
    \end{tikzpicture}%
  }
}{}
\makeatother 
\coursecode{CS301}
\renewcommand{\thesection}{9.\arabic{section}}
\setlength{\emergencystretch}{3em}

\title{Coding Assignment 9: Merge, Quick, and Shell Sort}
\author{Auqib Hamid Lone\\[0.3em]
  \emph{Department of Computer Science \& Engineering}, \emph{GCET Ganderbal}}
\date{\today}

\begin{document}
\maketitle

\noindent\textbf{Due:} two weeks from issue. There is no automatic late penalty --- if you need more time, ask before the deadline and an extension will be granted (see the course policy in the syllabus).

\noindent This assignment autogrades Week~9's institute practicals \key{P9 (quick sort)}, \key{P10 (shell sort)}, and \key{P11 (merge sort)}. The contract is given to you in \texttt{sorting.h} --- \emph{do not modify it}. Copy \texttt{starter\_sorting.c} to \texttt{sorting.c} and fill in the five functions. The setting is the running calculator: its symbol table is an integer array, sorted once at startup and searched many times afterwards, so this module is the startup sort --- plus one measurement, the inversion count, that tells the profiler how far the table was from sorted.

\section{Why this problem}

Week~9's argument is that the elementary $\Theta(n^2)$ sorts are not wrong, only mispriced. Merge sort buys a guaranteed $\Theta(n\log n)$ with a linear buffer; quick sort buys in-place operation but must defend its pivot against sorted input; shell sort buys in-place simplicity with a gap-dependent bound. The lecture builds all three, proves that no comparison sort can beat $\Omega(n\log n)$, and shows two ways around the wall (heap sort's $O(1)$ space and radix sort's changed model). This assignment is the three comparison sorts as running code. The one extension beyond the deck's listings is \texttt{countInversions}, which uses the merge step to count out-of-order pairs in $O(n\log n)$ --- a measurement that a naive double loop would take $\Theta(n^2)$.

\section{Your task}

Implement the five functions declared in \texttt{sorting.h}, all operating on \texttt{A[0..n-1]}, \texttt{n >= 0}:

\begin{verbatim}
void mergeRuns(int *A, int lo, int mid, int hi, int *tmp);
void mergeSort(int *A, int n);
void quickSort(int *A, int n);
void shellSort(int *A, int n);
long long countInversions(const int *A, int n);
\end{verbatim}

\textbf{Merge sort (P11).} \texttt{mergeRuns} merges two adjacent sorted runs using the caller's buffer \texttt{tmp}; on equal keys, take the \emph{left} run first (the stable merge). \texttt{mergeSort} allocates one buffer for the whole sort, recurses on the two halves, merges on the way back up, and frees the buffer. It must be $\Theta(n\log n)$ in every case and stable.

\textbf{Quick sort (P9).} \texttt{quickSort} is in-place and must not degrade to quadratic time on the inputs a symbol table actually produces: \emph{sorted}, \emph{reverse-sorted}, and \emph{all-equal}. Use median-of-three or a randomized pivot, and bound the recursion stack to $O(\log n)$ by recursing on the smaller side first. A three-way partition is one clean way to survive the all-equal case.

\textbf{Shell sort (P10).} \texttt{shellSort} is in-place with $O(1)$ auxiliary space and uses a decreasing gap sequence ending at \texttt{gap = 1}. Shell's halving sequence is acceptable; Knuth's $(3^k-1)/2$ sequence (1, 4, 13, 40, \dots) is better. The final gap-1 pass must be ordinary insertion sort.

\textbf{Inversion count.} \texttt{countInversions} returns the number of pairs $i<j$ with $A[i]>A[j]$, in $O(n\log n)$ worst case, \emph{without modifying} \texttt{A}. Use the merge step: when the right run's element is taken before the left run is exhausted, every remaining left element forms an inversion with it. Return a \texttt{long long} --- $n$ can reach $200{,}000$, where $n(n-1)/2$ exceeds \texttt{int}.

\section{Constraints}

\begin{itemize}
  \item Compile with \texttt{gcc -std=c11 -Wall -Wextra -Werror}: no errors and no warnings. No global or static mutable state.
  \item \texttt{n == 0} is the empty array (\texttt{A} may be \texttt{NULL}); every function must handle it. \texttt{n == 1} leaves the array unchanged.
  \item \texttt{mergeSort}: $\Theta(n\log n)$ worst case, $\Theta(n)$ auxiliary, stable, at most a constant number of allocations.
  \item \texttt{quickSort}: no $O(n)$ allocation; $O(\log n)$ stack; must finish on $200{,}000$ sorted, reverse-sorted, and all-equal elements within the autograder's timeout.
  \item \texttt{shellSort}: no \texttt{malloc}; gap sequence ends at $1$.
  \item \texttt{countInversions}: $O(n\log n)$; treats \texttt{A} as read-only (copy it if you need to sort); returns \texttt{long long}.
  \item Every array is sorted in \emph{ascending} (nondecreasing) order. Duplicates and negative values are legal.
\end{itemize}

\section{Worked examples (these are the visible tests)}

\textbf{Example 1 --- the merge step.} On \texttt{A = [1, 4, 7, 2, 5, 8]} with \texttt{lo = 0}, \texttt{mid = 2}, \texttt{hi = 5}, \texttt{mergeRuns} produces \texttt{[1, 2, 4, 5, 7, 8]}.

\textbf{Example 2 --- merge sort.} \texttt{mergeSort([5, 1, 4, 2, 8, 0, 3])} produces \texttt{[0, 1, 2, 3, 4, 5, 8]}; a single element is unchanged.

\textbf{Example 3 --- quick sort on adversarial input.} \texttt{quickSort} sorts \texttt{[1,2,3,4,5,6,7,8]}, \texttt{[8,7,6,5,4,3,2,1]}, \texttt{[4,4,4,4,4]}, and \texttt{[3,1,2,1,3,2]} to \texttt{[1,1,2,2,3,3]} for the last. A last-element Lomuto pivot passes the random cases but fails the sorted case's timeout --- that is the point of the constraint.

\textbf{Example 4 --- shell sort.}
\begin{verbatim}
shellSort([9,8,7,6,5,4,3,2,1]) -> [1,2,3,4,5,6,7,8,9]
shellSort([3,1,2,1,3,2])       -> [1,1,2,2,3,3]
\end{verbatim}

\textbf{Example 5 --- inversion count.} \texttt{countInversions([2,4,1,3])} is $3$: the pairs are $(2,1)$, $(4,1)$, $(4,3)$. \texttt{[1,2,3]} gives $0$; \texttt{[3,2,1]} gives $3$; \texttt{[2,2,2]} gives $0$ (equal values are not inversions); \texttt{[5,4,3,2,1]} gives $10 = 5\cdot4/2$.

\section{How to work}

\begin{enumerate}
  \item Copy \texttt{starter\_sorting.c} to \texttt{sorting.c}; keep \texttt{sorting.h} unchanged.
  \item Build and run the visible tests:
\begin{verbatim}
gcc -std=c11 -Wall -Wextra -Werror -o test_visible test_visible.c sorting.c
./test_visible
\end{verbatim}
  \item Trace one merge, one partition, and one gap pass on paper before coding --- the deck's figures are the model.
  \item Run the hidden suite only through the autograder; it adds random, sorted, reverse, all-equal, duplicate-heavy, organ-pipe, and $200{,}000$-element cases.
\end{enumerate}

\section{Submission and grading}

Submit \texttt{sorting.c} through the course submission form. Grading is by the autograder (correctness on the hidden suite) plus a short code read for the three complexity requirements: merge sort's $\Theta(n\log n)$ every case, quick sort's bounded stack on adversarial input, and the inversion count's merge-based $O(n\log n)$ with an unmodified input. A working sort that violates a stated complexity is partially credited; a warning under \texttt{-Werror} is a compile failure.

\bibliography{../../Bibliography_base}
\bibliographystyle{plain}
\end{document}
